% =====================================================================
% --- CAPITOLO 8: VERIFICA E VALIDAZIONE ---
% =====================================================================
\chapter{Verifica e Validazione}
\label{cap:testing}
% Riferimento Traccia: "Output: Testing e valutazione dell’usabilità"

Questo capitolo descrive le strategie di verifica (statica e dinamica) e validazione adottate per garantire la qualità del software DietiEstates25.

\section{Verifica Statica (Analisi del Codice)}
Per garantire la manutenibilità e prevenire difetti (\textit{Fault Avoidance}), è stata adottata una strategia di \textbf{Verifica Statica} continua.
Il progetto utilizza \textbf{ESLint} (versione 9.x) con la configurazione rigorosa \texttt{@antfu/eslint-config}.

Le principali regole di qualità applicate includono:
\begin{itemize}
    \item \textbf{Stile del Codice:} Indentazione a 2 spazi, punto e virgola obbligatorio e virgolette doppie, per garantire uniformità nel team.
    \item \textbf{Code Smells \& Best Practices:} La regola \texttt{no-console} (impostata a \textit{warn}) scoraggia l'uso di log di debug in produzione, prevenendo "rumore" nei log e potenziali problemi di sicurezza.
    \item \textbf{Convenzioni di Naming:} La regola \texttt{unicorn/filename-case} impone lo stile \textit{kebab-case} per tutti i file (es. \texttt{appointments.router.ts}), facilitando la navigazione nel progetto.
    \item \textbf{Import Sorting:} La regola \texttt{perfectionist/sort-imports} mantiene gli import ordinati automaticamente, riducendo drasticamente i conflitti di merge su file modificati da più sviluppatori.
\end{itemize}
L'analisi viene eseguita localmente tramite lo script \texttt{npm run lint} e validata tramite \texttt{npm run type-check} per assicurare la coerenza dei tipi TypeScript prima di ogni commit.

\section{Verifica Dinamica (Unit Testing)}
% Riferimento Slide: "Unit Testing in Isolation" e "Testing BB e WB"
La strategia di verifica dinamica si è concentrata sul \textbf{Unit Testing} dei componenti critici del backend. È stato adottato un approccio \textbf{Black-Box Testing}, basato sulla partizione in \textbf{Classi di Equivalenza} (Equivalence Class Partitioning) per definire gli scenari di test.

Per garantire l'isolamento dei test (\textit{Test in Isolation}), tutte le dipendenze esterne (come il database e i servizi di terze parti) sono state sostituite da \textbf{Mock Objects} utilizzando le funzionalità di \texttt{jest.mock} e \texttt{jest.spyOn}.

Di seguito sono riportati i dettagli di quattro metodi significativi sottoposti a test.

% --- METODO 1 ---
\subsection{Test del metodo \texttt{deleteAppointment}}

\begin{description}
    \item[Descrizione] Il metodo \texttt{deleteAppointment} (in \texttt{appointments.repository.ts}) elimina un appuntamento dal database, applicando filtri di sicurezza in base al ruolo dell'utente (es. un cliente può cancellare solo i propri appuntamenti).
    
    \item[Strategia di Testing] Approccio Black-Box con criterio R-WECT (Robust Weak Equivalence Class Testing). Il test verifica che la query al database venga costruita correttamente in base al ruolo, senza eseguire realmente la cancellazione (Isolamento tramite Mock di Drizzle).
    
    \item[Classi di Equivalenza Individuate]
    \begin{itemize}
        \item \textbf{CE1 (Valida - Admin):} Input: \texttt{id=1, userRole="admin"}. $\rightarrow$ Atteso: Query eseguita senza filtri aggiuntivi.
        \item \textbf{CE2 (Valida - Customer):} Input: \texttt{id=1, userId=123, userRole="customer"}. $\rightarrow$ Atteso: Query con clausola \texttt{WHERE userId = 123}.
        \item \textbf{CE3 (Non Valida):} Input: \texttt{id} inesistente. $\rightarrow$ Atteso: Nessuna riga eliminata (gestito dal DB).
    \end{itemize}
\end{description}

\begin{lstlisting}[language=Java, caption={Unit Test per deleteAppointment con Mocking}]
it("dovrebbe applicare filtri aggiuntivi se l'utente e' customer", async () => {
  // Arrange: Mock della catena delete -> where -> returning
  const { whereMock } = mockDeleteChain([{ id: 1 }]);

  // Act
  await deleteAppointment(1, 123, "customer");

  // Assert: Verifica isolamento e logica
  expect(db.delete).toHaveBeenCalled();
  expect(whereMock).toHaveBeenCalled();
});
\end{lstlisting}

% --- METODO 2 ---
\subsection{Test del metodo \texttt{getPostsNearby}}

\begin{description}
    \item[Descrizione] Il metodo \texttt{getPostsNearby} (in \texttt{posts.repository.ts}) esegue una ricerca geospaziale per trovare immobili entro un raggio specificato.
    
    \item[Strategia di Testing] Verifica della corretta costruzione della query spaziale. La funzione \texttt{db.select} e i metodi concatenati (\texttt{from}, \texttt{where}) sono stati "stubbati" per restituire dati fittizi.
    
    \item[Classi di Equivalenza Individuate]
    \begin{itemize}
        \item \textbf{CE1 (Valida):} Input: \texttt{lat=40.85, lon=14.26, radius=5000}. $\rightarrow$ Atteso: Lista immobili ordinati per distanza.
        \item \textbf{CE2 (Limite):} Input: \texttt{radius=0}. $\rightarrow$ Atteso: Lista vuota o query valida senza risultati.
    \end{itemize}
\end{description}

\begin{lstlisting}[language=Java, caption={Unit Test per getPostsNearby}]
it("dovrebbe eseguire la query spaziale con i parametri corretti", async () => {
  // Arrange: Mock della catena fluente di Drizzle
  const whereMock = jest.fn().mockReturnValue({ orderBy: orderByMock });
  const fromMock = jest.fn().mockReturnValue({ where: whereMock });
  (db.select as jest.Mock).mockReturnValue({ from: fromMock });

  // Act
  const result = await getPostsNearby(40.85, 14.26, 5000, 1, 10);

  // Assert
  expect(fromMock).toHaveBeenCalled(); 
  expect(result).toEqual(mockPosts);
});
\end{lstlisting}

% --- METODO 3 ---
\subsection{Test del metodo \texttt{getAllAppointmentsForCalendar}}

\begin{description}
    \item[Descrizione] Il metodo recupera gli appuntamenti formattati per il calendario, costruendo dinamicamente la query in base ai filtri di data (\texttt{start}, \texttt{end}).
    
    \item[Strategia di Testing] Test strutturale della logica condizionale di costruzione della query.
    
    \item[Classi di Equivalenza Individuate]
    \begin{itemize}
        \item \textbf{CE1 (Valida - No Filtri):} Input: \texttt{start=undefined}. $\rightarrow$ Atteso: Query senza clausole temporali.
        \item \textbf{CE2 (Valida - Range):} Input: \texttt{start=2025-06-01, end=2025-06-30}. $\rightarrow$ Atteso: Query con clausole \texttt{WHERE} su date.
        \item \textbf{CE3 (Non Valida - Logica):} Input: \texttt{start > end}. $\rightarrow$ Atteso: Risultato vuoto (query semanticamente valida ma logicamente impossibile).
    \end{itemize}
\end{description}

% --- METODO 4 ---
\subsection{Test dell'Endpoint \texttt{POST /appointments}}

\begin{description}
    \item[Descrizione] Questo metodo nel Router gestisce la creazione di un appuntamento, orchestrando la validazione dell'esistenza di immobile e agente prima di salvare.
    
    \item[Strategia di Testing] Test di integrazione dei componenti (Controller) con Mock dei Repository (\texttt{postsRepository}, \texttt{userRepository}). Le asserzioni (\texttt{expect}) fungono da \textbf{Oracolo} per verificare i codici di stato HTTP.
    
    \item[Classi di Equivalenza Individuate]
    \begin{itemize}
        \item \textbf{CE1 (Valida):} Input: \texttt{postId} esistente, \texttt{agentId} valido. $\rightarrow$ Atteso: \texttt{201 Created}, Email inviata.
        \item \textbf{CE2 (Non Valida - Immobile):} Input: \texttt{postId} inesistente. $\rightarrow$ Atteso: \texttt{404 Not Found} ("Immobile non trovato").
        \item \textbf{CE3 (Non Valida - Agente):} Input: \texttt{postId} valido ma agente inesistente. $\rightarrow$ Atteso: \texttt{404 Not Found} ("Agente non trovato").
    \end{itemize}
\end{description}

\begin{lstlisting}[language=Java, caption={Logica di validazione nel Router}]
// Logica testata tramite Mock dei Repository
const post = await postsRepository.getPostById(req.body.postId);
if (!post) {
  res.status(404).json({ message: "Immobile non trovato" });
  return;
}
// ...
\end{lstlisting}

\section{Valutazione dell'Usabilità}
% Riferimento Traccia: "Descrizione delle tecniche di valutazione dell’usabilità applicate"
% Riferimenti Slide: 40 - UI Validation.pdf

La valutazione dell'usabilità è stata condotta adottando un approccio misto, combinando metodi ispettivi (Expert Review) con test empirici su utenti reali, al fine di identificare problemi di interfaccia e migliorare l'esperienza utente (UX).

\subsection{Expert Reviews (Ispezione Euristica)}
% Riferimento Traccia: "definizione di una checklist per review sistematica..."
È stata condotta un'ispezione esperta basata sulle \textbf{8 Regole d'Oro di Shneiderman}. L'analisi ha evidenziato una buona aderenza ai principi di design, con alcune criticità minori identificate e risolte (o tracciate per sviluppi futuri).

\begin{longtable}{|p{8cm}|c|p{6cm}|}
\hline
\textbf{Euristica di Shneiderman} & \textbf{Esito} & \textbf{Note / Osservazioni} \\ \hline
\endfirsthead
\multicolumn{3}{c}{{\bfseries \tablename\ \thetable{} -- segue dalla pagina precedente}} \\ \hline
\textbf{Euristica} & \textbf{Esito} & \textbf{Note} \\ \hline
\endhead
\hline \multicolumn{3}{r}{{\emph{continua nella pagina seguente}}} \\ \endfoot
\hline
\endlastfoot

1. Sforzarsi per la coerenza & \checkmark & Layout e colori uniformi tra Home, Ricerca e Dettaglio. \\ \hline
2. Abilitare scorciatoie per utenti esperti & $\sim$ & Mancano hotkeys da tastiera per la ricerca rapida. \\ \hline
3. Fornire un feedback informativo & \checkmark & Toast notifications (ngx-toastr) presenti per ogni azione (login, errore, salvataggio). \\ \hline
4. Progettare dialoghi di chiusura & \checkmark & I form (es. Inserimento Annuncio) mostrano conferma prima dell'invio. \\ \hline
5. Prevenire gli errori & \checkmark & Validazione input in tempo reale (Reactive Forms) impedisce l'invio di dati errati. \\ \hline
6. Consentire annullamenti semplici & \textbf{NO} & Manca un tasto "Undo" immediato dopo l'eliminazione di un preferito. \\ \hline
7. Favorire il controllo dell'utente & \checkmark & L'utente può navigare liberamente e resettare i filtri di ricerca. \\ \hline
8. Ridurre il carico cognitivo & \checkmark & Interfaccia minimalista, informazioni divise in "card" leggibili. \\
\hline
\caption{Checklist basata sulle 8 Regole d'Oro di Shneiderman}
\label{tab:expert_review}
\end{longtable}

\subsection{Esperimento con Utenti Reali}
% Riferimento Traccia: "Progettazione e conduzione di un esperimento con utenti reali..."

Per validare l'usabilità del sistema in uno scenario reale, è stato condotto un esperimento controllato.

\subsubsection{Descrizione dei Soggetti e Procedura}
\begin{description}
    \item[Campione Utenti:] Sono stati reclutati \textbf{5 partecipanti} non coinvolti nello sviluppo:
    \begin{itemize}
        \item 3 studenti universitari fuori sede (target "Cliente Giovane").
        \item 2 lavoratori (target "Cliente Adulto").
    \end{itemize}
    \item[Procedura:] Ad ogni utente è stato chiesto di eseguire i seguenti \textbf{3 Task} senza ricevere aiuto esterno, utilizzando l'applicazione su browser Desktop (Chrome):
    \begin{enumerate}
        \item \textbf{Task A (Ricerca):} Cercare un appartamento a "Napoli" con prezzo max 600€.
        \item \textbf{Task B (Prenotazione):} Prenotare una visita per il primo risultato trovato, verificando il meteo.
        \item \textbf{Task C (Login/Dashboard):} Effettuare il login e visualizzare la lista dei propri appuntamenti.
    \end{enumerate}
\end{description}

\subsubsection{Metriche Rilevate}
Durante l'esperimento sono state misurate:
\begin{itemize}
    \item \textbf{Efficacia (Success Rate):} Percentuale di task completati senza errori critici.
    \item \textbf{Efficienza (Time on Task):} Tempo medio impiegato per completare il task.
    \item \textbf{Error Rate:} Numero medio di click errati o percorsi sbagliati per utente.
\end{itemize}

\begin{table}[h!]
    \centering
    \begin{tabular}{|l|c|c|c|}
    \hline
    \textbf{Task} & \textbf{Success Rate} & \textbf{Tempo Medio (sec)} & \textbf{Errori Medi} \\ \hline
    Task A (Ricerca) & 100\% & 45s & 0.2 \\ \hline
    Task B (Prenotazione) & 80\% & 90s & 1.5 \\ \hline
    Task C (Login) & 100\% & 20s & 0.0 \\ \hline
    \end{tabular}
    \caption{Risultati metrici dell'esperimento di usabilità.}
    \label{tab:usability_metrics}
\end{table}

\subsubsection{Survey Post-Test e Soddisfazione}
Al termine della sessione, agli utenti è stato somministrato un questionario (basato su scala Likert 1-5, dove 5 = Molto Soddisfatto) per valutare la soddisfazione soggettiva.

\begin{description}
    \item[Risultati del Survey:]
    \begin{enumerate}
        \item \textit{"È stato facile trovare l'immobile?"} $\rightarrow$ \textbf{Media: 4.8/5}
        \item \textit{"L'interfaccia è chiara e leggibile?"} $\rightarrow$ \textbf{Media: 4.6/5}
        \item \textit{"Il processo di prenotazione è stato fluido?"} $\rightarrow$ \textbf{Media: 3.8/5}
        \item \textit{"Le previsioni meteo sono state utili?"} $\rightarrow$ \textbf{Media: 4.2/5}
    \end{enumerate}
\end{description}

\subsubsection{Discussione dei Risultati}
L'esperimento ha dimostrato un'elevata usabilità generale del sistema DietiEstates25.
\begin{itemize}
    \item \textbf{Punti di Forza:} La funzione di ricerca (Task A) è risultata estremamente intuitiva e veloce, con un tasso di successo del 100\%. Gli utenti hanno apprezzato particolarmente la pulizia dell'interfaccia Angular.
    \item \textbf{Criticità:} Nel Task B (Prenotazione), un utente ha avuto difficoltà a selezionare l'orario nel calendario, causando un tasso di errore più alto. Questo suggerisce la necessità di migliorare il componente \textit{Datepicker} nelle versioni future.
    \item \textbf{Conclusione:} Con un punteggio medio di soddisfazione di \textbf{4.35/5}, il sistema è considerato pronto per il rilascio, previa correzione minore sul calendario prenotazioni.
\end{itemize}